%%%%%%%%%%%%%%%%%%%%%%%%%%%%%%%%%%%%%%%%%
% Short Sectioned Assignment LaTeX Template Version 1.0 (5/5/12)
% This template has been downloaded from: http://www.LaTeXTemplates.com
% Original author:  Frits Wenneker (http://www.howtotex.com)
% License: CC BY-NC-SA 3.0 (http://creativecommons.org/licenses/by-nc-sa/3.0/)
%%%%%%%%%%%%%%%%%%%%%%%%%%%%%%%%%%%%%%%%%

% \documentclass[paper=a4, fontsize=11pt]{scrartcl} % A4 paper and 11pt font size
\documentclass[11pt, a4paper]{book}
\usepackage[T1]{fontenc} % Use 8-bit encoding that has 256 glyphs
\usepackage[utf8]{inputenc}
\usepackage{fourier} % Use the Adobe Utopia font for the document - comment this line to return to the LaTeX default
\usepackage{listings} % para insertar código con formato similar al editor
\usepackage[spanish, es-tabla]{babel} % Selecciona el español para palabras introducidas automáticamente, p.ej. "septiembre" en la fecha y especifica que se use la palabra Tabla en vez de Cuadro
\usepackage{url} % ,href} %para incluir URLs e hipervínculos dentro del texto (aunque hay que instalar href)
\usepackage{graphics,graphicx, float} %para incluir imágenes y colocarlas
\usepackage[gen]{eurosym} %para incluir el símbolo del euro
\usepackage{cite} %para incluir citas del archivo <nombre>.bib
\usepackage{enumerate}
\usepackage{hyperref}
\usepackage{graphicx}
\usepackage{tabularx}
\usepackage{booktabs}

\usepackage[table,xcdraw]{xcolor}
\hypersetup{
	colorlinks=true,	% false: boxed links; true: colored links
	linkcolor=black,	% color of internal links
	urlcolor=cyan		% color of external links
}
\renewcommand{\familydefault}{\sfdefault}
\usepackage{fancyhdr} % Custom headers and footers
\pagestyle{fancyplain} % Makes all pages in the document conform to the custom headers and footers
\fancyhead[L]{} % Empty left header
\fancyhead[C]{} % Empty center header
\fancyhead[R]{Me Llamo Así} % My name
\fancyfoot[L]{} % Empty left footer
\fancyfoot[C]{} % Empty center footer
\fancyfoot[R]{\thepage} % Page numbering for right footer
%\renewcommand{\headrulewidth}{0pt} % Remove header underlines
\renewcommand{\footrulewidth}{0pt} % Remove footer underlines
\setlength{\headheight}{13.6pt} % Customize the height of the header

\usepackage{titlesec, blindtext, color}
\definecolor{gray75}{gray}{0.75}
\newcommand{\hsp}{\hspace{20pt}}
\titleformat{\chapter}[hang]{\Huge\bfseries}{\thechapter\hsp\textcolor{gray75}{|}\hsp}{0pt}{\Huge\bfseries}
\setcounter{secnumdepth}{4}
\usepackage[Lenny]{fncychap}


\begin{document}

	% Plantilla portada UGR
	\begin{titlepage}
\newlength{\centeroffset}
\setlength{\centeroffset}{-0.5\oddsidemargin}
\addtolength{\centeroffset}{0.5\evensidemargin}
\thispagestyle{empty}

\noindent\hspace*{\centeroffset}\begin{minipage}{\textwidth}

\centering
\includegraphics[width=0.9\textwidth]{logos/logo_ugr.jpg}\\[1.4cm]

\textsc{ \Large TRABAJO FIN DE GRADO\\[0.2cm]}
\textsc{ GRADO EN INGENIERÍA INFORMÁTICA}\\[1cm]

{\Huge\bfseries Programación eficiente de un algoritmo de procesamiento de la actividad cerebral \\}
\noindent\rule[-1ex]{\textwidth}{3pt}\\[3.5ex]
{\large\bfseries Implementación y optimización mediante paralelización de un algoritmo de análisis fractal de matrices binarias }
\end{minipage}

\vspace{1cm}
\noindent\hspace*{\centeroffset}
\begin{minipage}{\textwidth}
\centering

\textbf{Autor}\\ {Miguel Ángel Posadas Arráez}\\[2.5ex]
\textbf{Director}\\ Juan Ruiz de Miras\\[2cm]
\includegraphics[width=0.3\textwidth]{logos/etsiit_logo.png}\\[0.1cm]
\textsc{Escuela Técnica Superior de Ingenierías Informática y de Telecomunicación}\\
\textsc{---}\\
Granada, Junio de 201x
\end{minipage}
\end{titlepage}


	% Plantilla prefacio UGR
	\thispagestyle{empty}

\begin{center}
{\large\bfseries Programación eficiente de un algoritmo de procesamiento de la actividad cerebral \\ Implementación y optimización mediante paralelización de un algoritmo de análisis fractal de matrices binarias }\\
\end{center}
\begin{center}
Miguel Ángel Posadas Arráez\\
\end{center}

%\vspace{0.7cm}

\vspace{0.5cm}
\noindent{\textbf{Palabras clave}: \textit{Box counting, análisis fractal de matrices binarias, programación paralela, programación GPU}
\vspace{0.7cm}

\noindent{\textbf{Resumen}\\
El procesamiento de datos es una parte esencial en el estudio de diversas enfermedades neurodegenerativas. No obstante, las técnicas actuales para la adquisión de estos datos(resonancia magnética, encefalografía etc...) generan grandes volumenes de datos, lo que conlleva que su procesamiento tenga un gran coste computacional. El propósito de este trabajo es la programación eficiente del método conocido como Box counting para el procesamiento de matrices binarias.  Para la toma de tiempos y con el objetivo de poder hacer diversas comparaciones hemos utilizado dos plataformas con software y hardware diferentes. En un principio, el algoritmo venía programado en Matlab, con el objetivo de obtener una versión más eficiente de este algoritmo implementé una versión secuencial en C++ y posteriormente traté de aprovechar el paralelismo dado por los procesadores multi-core introduciendo la tecnología OpenACC obteniendo aceleraciones de hasta 15.7x. Finalmente, para la obtención de mejores aceleraciones, utilizamos el paradigma de computación de propósito general en unidades de procesamiento gráfico (GPGPU) para aprovechar la arquitectura \textit{Single Instruction, Multiple Data} (SIMD) con el que obtenimos aceleraciones de hasta 77.47x .
	

\cleardoublepage

\begin{center}
	{\large\bfseries Efficient programming of a cerebral processing activity algorithm \\ Implemantation and optimization with parallelization of a binary matrixes fractal analysis algorithm }\\
\end{center}
\begin{center}
	Miguel Ángel Posadas Arráez\\
\end{center}
\vspace{0.5cm}
\noindent{\textbf{Keywords}: \textit{Box counting, Fractal dimension of binary matrix, parallel programming, GPU programming}
\vspace{0.7cm}

\noindent{\textbf{Abstract}\\
Data processing is essential in neurodegenerative diseases studies. However, the existing ways of getting this data (magnetic resonance, electro-encephalography ... ) generate big data volumes, so processing that data has a expensive computing cost. The purpose of this paper is the efficient programming of the Box counting method in order to apply it at binary matrixes processing. They have been used two differents hardware and software platforms to track the time and making comparisons between then. At first, the algorithm was write on Matlab. I \textit{"translated"} that code to C++ with the goal of getting a more efficient sequential algorithm. Then, I tried to parallelize it with OpenACC with the goal of using the multicore CPU and I get a 15.7x speedup. Finally, in order to get the best speedup, I used the general-purpose computing on graphics processing units (GPGPU) to take advantage of the \textit{Single Instruction, Multiple Data} (SIMD) architecture on the graphics processing unit (GPU), getting a speedup of 77.47x.



\cleardoublepage

\thispagestyle{empty}

\noindent\rule[-1ex]{\textwidth}{2pt}\\[4.5ex]

D. \textbf{Juan Ruiz de Miras}, Profesor del Departamento de Lenguajes y Sistemas Informáticos

\vspace{0.5cm}

\textbf{Informo:}

\vspace{0.5cm}

Que el presente trabajo, titulado \textit{\textbf{Programación eficiente de un algoritmo de procesamiento de la actividad cerebral}},
ha sido realizado bajo mi supervisión por \textbf{Miguel Ángel Posadas Arráez}, y autorizo la defensa de dicho trabajo ante el tribunal
que corresponda.

\vspace{0.5cm}

Y para que conste, expiden y firman el presente informe en Granada a Julio de 2021.

\vspace{1cm}

\textbf{El/la director(a)/es: }

\vspace{5cm}

\noindent \textbf{Juan Ruiz de Miras}

\chapter*{Agradecimientos}






	% Índice de contenidos
	\newpage
	\tableofcontents

	% Índice de imágenes y tablas
	\newpage
	\listoffigures

	% Si hay suficientes se incluirá dicho índice
	\listoftables 
	\newpage

	% Introducción 
	\chapter{Introducción}


El cálculo de la dimensión fractal de un conjunto de datos es una técnica con una gran multitud de aplicaciones, desde el ámbito económico, -- incluyendo su relación con la premisa del Bitcoin --  {\cite{delfin2016fractal}}, hasta el campo de la medicina y estudio del cuerpo humano, siendo este cálculo útil en la detección de tumores cerebrales {\cite{iftekharuddin2009fractal}}.

La dimensión fractal tiene su cabida en la neurociencia, siendo útil para el procesamiento de la actividad cerebral con el fin de poder estudiar diversas patologías neurodegenerativas \cite{fernandez2001use}. Sin embargo las técnicas actuales para la adquisición de este tipo de datos tales como las electroencefalografías o las resonancias magnéticas, generan grandes volúmenes de datos que implican un alto coste computacional, ya que en función de la dimensión del conjunto de datos, el orden del algoritmo aumenta de manera exponencial, siendo O(n) para conjuntos de datos de una dimensión y O($n^{4}$)  para matrices de cuatro dimensiones.

El propósito de este proyecto es el de la implementación eficiente del algoritmo conocido como Box counting. Para ello, partiendo del algoritmo, inicialmente escrito en el lenguaje de programación Matlab se desarrolla una versión secuencial, escrita en C++ que sirve como base para utilizar diversas tecnologías,como OpenACC y CUDA, con el fin de aprovechar el paralelismo que proporcionan los microprocesadores actuales así como la arquitectura SIMD (\textit{Single Instruction Multiple Data}) presente en las tarjetas gráficas.

	% Descripción del problema y hasta donde se llega
	\chapter{Descripción del problema}



	% Estado del arte
	% 	1. Crítica al estado del arte
	% 	2. Propuesta
	\chapter{Estado del arte}

El software libre y sus licencias \cite{gplv3} ha permitido llevar a cabo una expansión del
aprendizaje de la informática sin precedentes.

	
	\chapter{Planificación}

\section{Metodología utilizada}


\section{Temporización}

\section{Seguimiento del desarrollo}


	% Análisis del problema
	% 1. Análisis de requisitos
	% 2. Análisis de las soluciones
	% 3. Solucion propuesta
	% 4. Análisis de seguridad
	\chapter{Análisis del problema}
 


	% Desarrollo bajo sprints: 
	% 	1. Permitir registros y login de usuarios
	% 	2. Desarrollo del sistema de incidencias
	% 	3. Desarrollo del sistema de denuncias administrativas y accidentes
	% 	4. Desarrollo del sistema de croquis
	%   5. Instalación de la aplicación de manera automática
	\chapter{Implementación}

La implementación del software se ha dividido en hitos. Estos, han sido definidos en Github
y cada uno de ellos contiene un grupo de \textit{issues} que se corresponden con las distintas
mejoras que se han ido incorporando al software a lo largo de su desarrollo.\\



	% Presupuesto

	% Conclusiones
	\chapter{Conclusiones y trabajos futuros}



	% Trabajos futuros


	
	\newpage
	\bibliography{bibliografia}
	\bibliographystyle{plain}
	
\end{document}

